We investigate the loss landscape geometry of neural theorem provers under two
proof generation strategies: chain-of-thought (CoT) decomposition, which
produces step-by-step derivations, and direct generation, which outputs
conclusions in a single step. Using small transformer models trained on
propositional logic theorem proving with matched architectures and training
conditions, we measure five geometric properties of the loss landscape: Hessian
top eigenvalue, Hessian trace, sharpness-aware minimization (SAM) sharpness,
mode connectivity barriers, and perturbation sensitivity. We prove that under a
conditional independence assumption on intermediate reasoning steps, the
spectral norm of the Hessian of the decomposed loss is bounded by the maximum
per-step Hessian spectral norm, which is generically smaller than the spectral
norm of the monolithic loss. Our experiments confirm that CoT models converge to
flatter minima (top eigenvalue $3.20 \pm 0.78$ vs.\ $4.16 \pm 0.47$), exhibit
lower SAM sharpness at all perturbation radii, and maintain smaller gradient
norms during training. Unexpectedly, CoT solutions display \emph{higher} loss
barriers under linear interpolation between independently trained models ($1.40$
vs.\ $1.07$), indicating that while CoT minima are individually flatter, they
are more isolated in parameter space. These results provide the first geometric
characterization of how proof decomposition reshapes the optimization landscape
in neural theorem proving.
